\documentclass[11pt,a4paper]{ctexart}

\usepackage[margin=2.5cm]{geometry}
\usepackage{amsmath,amssymb}
\usepackage{booktabs}
\usepackage{tabularx}
\usepackage{enumitem}
\usepackage{xcolor}
\usepackage{hyperref}

\hypersetup{
    colorlinks=true,
    linkcolor=blue!60!black,
    urlcolor=blue!60!black,
    citecolor=blue!60!black
}

\setlist[itemize]{leftmargin=2em,itemsep=0.25em,topsep=0.25em}
\setlist[enumerate]{leftmargin=2em,itemsep=0.25em,topsep=0.25em}

\newcommand{\eps}{\varepsilon}
\newcommand{\dd}{\,\mathrm d}
\newcommand{\argmax}{\operatorname*{arg\,max}}
\newcommand{\todoitem}[1]{\item[$\square$] #1}

\title{WKB 稀有突变极限文章近期修改计划}
\author{}
\date{}

\begin{document}

\maketitle

\section{总体目标}

近期工作的目标是先完成文章非数值部分的润色和结构冻结。数值部分暂时不以得到全部结果为 deadline，而是先固定实验设计、图表位置、诊断量和结果说明方式。等数值结果出来后，只需要插入图表并补充相应讨论。

本文主线建议固定为
\begin{equation*}
\text{WKB 重构}
\quad \Longrightarrow \quad
\text{半离散格式设计}
\quad \Longrightarrow \quad
\text{结构性质与 stationary AP}
\quad \Longrightarrow \quad
\text{全离散演化策略}
\quad \Longrightarrow \quad
\text{数值验证}.
\end{equation*}

文章应突出以下四个亮点。
\begin{enumerate}
    \item 设计基于 WKB 重构的相位和振幅耦合半离散格式。
    \item 证明半离散层面的正性保持，并在固定 trait 网格下建立 stationary AP 结构。
    \item 指出全离散初值演化问题与稳态极限问题并不相同，并给出必要的时间离散策略。
    \item 通过数值实验验证 WKB 粗网格重构能够捕捉 direct density solver 在细 trait 网格下才能解析的 selected trait 和集中结构。
\end{enumerate}

\section{建议 deadline}

建议将非数值部分的硬 deadline 设为
\begin{center}
    \fbox{\large 6 月 1 日晚完成 clean TeX 和 clean PDF}
\end{center}

6 月 2 日上午只做小修，包括 typo、交叉引用、公式编号、图表占位和 bibliography 检查。原则上不再新增理论、不再重排结构、不再改变 AP 主张。

到 6 月 1 日晚应完成以下内容。
\begin{itemize}
    \item 全文结构调整完成。
    \item Introduction、理论部分、全离散说明、数值实验设计和 Conclusion 完成润色。
    \item 所有红色标注、临时说明和重复 remark 删除或改成正式表述。
    \item 数值部分保留图表占位和实验说明，后续只需要插图和补充结果讨论。
    \item 生成一版可编译的 clean TeX 和 clean PDF。
\end{itemize}

\section{建议文章结构}

建议将文章结构调整为以下形式。
\begin{verbatim}
\section{Introduction}

\section{WKB reformulation and continuous asymptotic target}

\section{Semi-discrete WKB scheme and stationary AP structure}

\section{Fully discrete realization for evolutionary computation}

\section{Numerical experiments}

\section{Conclusion}
\end{verbatim}

其中最重要的调整是将全离散部分从半离散分析中单独分出来。这样可以更清楚地区分 stationary semi-discrete AP structure 和 time-dependent fully discrete computation，也可以避免审稿人误解为本文已经证明了完整的 fully discrete AP theorem。

\section{各部分修改任务}

\subsection{Introduction}

Introduction 需要围绕数值困难和本文贡献重写。建议完成以下修改。
\begin{itemize}
    \item 说明 rare mutation regime 下原始密度 $n^\eps(x,\theta,t)$ 在 trait 方向高度集中。
    \item 强调 direct density solver 需要极细 trait 网格，因此计算成本高，并且容易受到尖峰分辨率的限制。
    \item 引入 WKB 重构
    \begin{equation*}
        n^\eps(x,\theta,t)
        = W^\eps(x,\theta,t)
        \exp\!\left(\frac{u^\eps(\theta,t)}{\eps}\right).
    \end{equation*}
    \item 将本文贡献明确写成 WKB 格式设计、半离散结构分析、stationary fixed-grid AP、全离散演化策略和数值验证。
    \item paper organization 保持简洁，不要写成过于机械的段落。
\end{itemize}

\subsection{WKB reformulation and continuous asymptotic target}

这一节只作为理论目标和数值设计动机，不应写得像本文主要证明连续极限。建议完成以下修改。
\begin{itemize}
    \item 检查原模型、WKB ansatz、相位方程和振幅方程的符号是否统一。
    \item 强调 $\max_\theta u^\eps$ 和 $\argmax_\theta u^\eps$ 用于识别 fittest trait。
    \item 连续 rare mutation limit 只作为 asymptotic target，不需要过度展开。
    \item 减少重复介绍，将重点留给半离散格式和 AP 结构分析。
\end{itemize}

\subsection{Semi-discrete WKB scheme and stationary AP structure}

这一节是文章理论核心，需要重点打磨。建议完成以下修改。
\begin{itemize}
    \item 统一半离散记号，包括 $x_j$、$\theta_k$、$\Delta x$、$\Delta\theta$、$W_{j,k}^\eps$、$u_k^\eps$、$E_k^\eps$、$n_{j,k}^\eps$ 和 $\rho_j^\eps$。
    \item 明确 WKB density reconstruction
    \begin{equation*}
        E_k^\eps=\exp\!\left(\frac{u_k^\eps}{\eps}\right),
        \qquad
        n_{j,k}^\eps = W_{j,k}^\eps E_k^\eps.
    \end{equation*}
    \item 半离散格式部分突出两个设计点，即数值 Hamiltonian 处理相位方程，conservative flux 处理振幅方程。
    \item 正性保持证明应尽量简洁，重点说明振幅方程具有 quasi-positive 或 Metzler-type 结构。
    \item weighted remainder 条件要重新组织，不要显得像临时补丁。应将其解释为 WKB 离散链式法则缺失导致的结构性 defect，并给出 one-sided stability assumption。
    \item AP 结论应统一表述为 stationary fixed-grid AP structure，不要写成完整的 time-dependent fully discrete AP theorem。
    \item density-compatible variant 可以保留，但应降低定位。它主要用于说明 weighted remainder 条件在某些结构下是合理的，而不是文章第二套主算法。
\end{itemize}

\subsection{Fully discrete realization for evolutionary computation}

这一节应单独成节。核心目标不是证明完整定理，而是说明如何可靠计算初值演化问题。建议完成以下修改。
\begin{itemize}
    \item 开头明确说明半离散 stationary AP 与全离散初值演化是两个不同问题。
    \item 解释 naive time discretization 的潜在问题。主要原因是 $u$ 的时间离散误差会通过 $\exp(u/\eps)$ 被放大。
    \item 写清楚本文采用或建议的时间离散策略，例如 $\eps$-aware time step、adaptive time step、semi-implicit phase update、projective strategy 或 positivity-preserving treatment。
    \item 不要将这一节写成未完成证明的 theorem。更合适的定位是 numerical realization 和 practical stability discussion。
    \item 保留后续数值实验需要的诊断量。
\end{itemize}

\subsection{Numerical experiments}

数值部分现在先完成实验设计和图表位置。结果出来后再补充具体讨论。建议按以下顺序组织。

\paragraph{Test 1. Structural diagnostics.}

该测试用于检验正性、密度重构、weighted remainder 和相位最大点等结构量是否表现合理。建议准备以下图表。
\begin{itemize}
    \item $\min W$ 随时间变化。
    \item weighted remainder diagnostic 随时间变化。
    \item $\max_\theta u^\eps$ 或 $\argmax_\theta u^\eps$ 随时间变化。
    \item density reconstruction 与总密度变化。
\end{itemize}

\paragraph{Test 2. Long-time trait concentration.}

该测试用于展示解在 trait 方向的集中，以及 $u^\eps$ 的最大点稳定指向优势 trait。建议准备以下图表。
\begin{itemize}
    \item 不同时刻的 trait marginal。
    \item 不同时刻的 $u^\eps(\theta,t)$ 曲线。
    \item $\argmax_\theta u^\eps(t)$ 随时间变化。
\end{itemize}

\paragraph{Test 3. Direct fine-grid solver versus WKB reconstruction.}

该测试用于比较 direct density solver 与 WKB reconstruction 的结果。建议比较以下量。
\begin{itemize}
    \item 同一时刻 direct 与 WKB 的 trait marginal。
    \item 空间总密度 $\rho(x)$。
    \item selected trait。
    \item 误差表。
\end{itemize}

\paragraph{Test 4. Trait-grid refinement of the direct density solver toward WKB coarse-grid observables.}

这是建议新增的重点实验。它用于说明 direct density solver 只有在 $\theta$ 网格足够细时，才能得到 WKB 粗网格重构已经能捕捉的结果。

建议设置为
\begin{equation*}
    K_{\rm WKB}=32 \quad \text{or} \quad 64,
\end{equation*}
以及
\begin{equation*}
    K_{\rm dir}=32,64,128,256,512.
\end{equation*}
如果计算量允许，可以加入 $K_{\rm dir}=1024$。所有计算固定同一组小 $\eps$、同一空间网格和同一终止时间。

WKB selected trait 定义为
\begin{equation*}
    \theta_{\rm WKB}=\argmax_k u_k^\eps.
\end{equation*}
Direct solver 的 selected trait 可定义为
\begin{equation*}
    \theta_{\rm dir}(K_{\rm dir})
    = \argmax_\theta P^{\rm dir}_{K_{\rm dir}}(\theta),
\end{equation*}
其中 $P^{\rm dir}_{K_{\rm dir}}$ 表示 direct density solver 得到的 trait marginal。

建议比较 selected trait 误差
\begin{equation*}
    e_\theta(K_{\rm dir})
    = d_{\mathbb T}(
    \theta_{\rm dir}(K_{\rm dir}),
    \theta_{\rm WKB}
    \bigr),
\end{equation*}
空间密度误差
\begin{equation*}
    e_\rho(K_{\rm dir})
    =
    \frac{
    \|\rho^{\rm dir}_{K_{\rm dir}}-\rho^{\rm WKB}\|_1
    }{
    \|\rho^{\rm WKB}\|_1
    },
\end{equation*}
以及 coarse-binned trait marginal 误差。具体而言，先将 direct 细网格 trait marginal 投影到 WKB 粗网格上，记为
\begin{equation*}
    \Pi_{K_{\rm WKB}} P^{\rm dir}_{K_{\rm dir}},
\end{equation*}
然后计算
\begin{equation*}
    e_P(K_{\rm dir})
    =
    \frac{
    \|\Pi_{K_{\rm WKB}} P^{\rm dir}_{K_{\rm dir}}
    - P^{\rm WKB}\|_1
    }{
    \|P^{\rm WKB}\|_1
    }.
\end{equation*}

建议准备以下图表。
\begin{itemize}
    \item $\theta_{\rm dir}(K_{\rm dir})$ 随 $K_{\rm dir}$ 变化的曲线，并以水平线标出 WKB 粗网格结果。
    \item direct trait marginal 随 $K_{\rm dir}$ 加密逐渐变尖，同时展示 WKB 粗网格结果已经定位峰值。
    \item $e_\theta$、$e_\rho$、$e_P$ 和 CPU time 随 $K_{\rm dir}$ 的变化。
\end{itemize}

该实验的结论建议写成如下形式。
\begin{quote}
Direct density solver 在 trait 网格加密后趋向于 WKB 粗网格重构所给出的 selected trait、空间密度和 coarse-binned trait marginal。相比之下，WKB 方法在较少 trait 网格点下已经能够稳定捕捉优势性状位置，因此具有明显的分辨率优势。
\end{quote}

需要注意的是，不建议声称 direct solver 的逐点尖峰密度完全趋向 WKB 粗网格密度。更稳妥的比较对象是 selected trait、空间总密度 $\rho(x)$ 和 coarse-binned trait marginal。

\paragraph{Test 5. Coarse trait-grid advantage.}

该测试用于展示 WKB 方法对 $K_\theta$ 不敏感，而 direct solver 对 $K_\theta$ 更敏感。它可以和 Test 4 合并，也可以单独作为效率实验。建议准备以下图表。
\begin{itemize}
    \item 不同 $K_\theta$ 下 WKB selected trait 的对比。
    \item 不同 $K_\theta$ 下 direct selected trait 的对比。
    \item CPU time 表。
    \item 误差和计算量表。
\end{itemize}

\subsection{Conclusion}

必须补一个正式结论。当前可以先写非数值版，后续等图出来后再补两三句数值结果总结。Conclusion 建议包括以下内容。
\begin{itemize}
    \item 本文提出了 WKB 半离散格式。
    \item 半离散格式具有正性保持性质，并在适当结构条件下给出 stationary fixed-grid AP 结构。
    \item 全离散演化问题与稳态极限问题不同，因此需要额外的时间离散策略。
    \item 数值实验将验证 WKB 粗网格方法能够捕捉 direct fine-grid 方法的极限行为。
    \item 后续可研究完整 fully discrete AP 理论、更高维问题、复杂生态模型和更强的重构策略。
\end{itemize}

\section{每日安排}

\begin{tabularx}{\textwidth}{@{}p{2.6cm}X X@{}}
\toprule
时间 & 主要任务 & 交付物 \\
\midrule
5 月 28 日 & 固定文章结构。重写 Introduction 的贡献段。检查 abstract 与贡献是否一致。不再增加新的理论目标。 & Introduction clean 版。Abstract 初步 clean 版。文章结构固定版。 \\
\addlinespace
5 月 29 日 & 修改 Section 2 和 Section 3 前半部分。将连续 WKB 部分写成 asymptotic target。统一 WKB ansatz 和 notation。检查数值 Hamiltonian 和 flux 假设。 & Section 2 clean 版。Section 3.1 到 3.2 clean 版。 \\
\addlinespace
5 月 30 日 & 修改半离散分析和 AP 结构。精简正性证明。重新组织 weighted remainder assumption。将 AP 结论统一为 stationary fixed-grid AP。 & 半离散分析 clean 版。Stationary AP theorem clean 版。AP scope remark clean 版。 \\
\addlinespace
5 月 31 日 & 单独整理 fully discrete 部分和数值实验骨架。写清楚初值演化与稳态极限的区别。加入 direct trait-grid refinement 实验设计。 & Fully discrete section clean 版。Numerical experiments skeleton clean 版。新增 Test 4 的完整文字说明。 \\
\addlinespace
6 月 1 日 & 全文通读和冻结非数值部分。删除红色文字。检查公式编号、定理编号和引用编号。补 Conclusion。编译 clean PDF。 & Clean TeX。Clean PDF。等待数值结果插入的冻结版。 \\
\addlinespace
6 月 2 日上午 & 最后一轮小检查，只处理 typo、交叉引用、图表占位、公式编号和 bibliography。 & 非数值部分最终版。 \\
\bottomrule
\end{tabularx}

\section{最终待办清单}

\subsection{必须完成}

\begin{itemize}
    \todoitem{重写 Introduction 贡献段。}
    \todoitem{修改 Section 2 标题和定位。}
    \todoitem{统一 WKB notation。}
    \todoitem{修改半离散格式描述。}
    \todoitem{精简正性证明。}
    \todoitem{重新组织 weighted remainder assumption。}
    \todoitem{将 AP 结论统一为 stationary fixed-grid AP。}
    \todoitem{单独写 fully discrete realization section。}
    \todoitem{明确初值演化与稳态极限不同。}
    \todoitem{加入 direct trait-grid refinement 实验设计。}
    \todoitem{补 Conclusion。}
    \todoitem{删除所有红色文字。}
    \todoitem{编译 clean PDF。}
\end{itemize}

\subsection{数值结果出来后完成}

\begin{itemize}
    \todoitem{插入结构诊断图。}
    \todoitem{插入 long-time concentration 图。}
    \todoitem{插入 direct fine-grid 与 WKB 对比图。}
    \todoitem{插入 direct trait-grid refinement 图。}
    \todoitem{插入 CPU time 和误差表。}
    \todoitem{根据图补充每个实验后的 discussion。}
    \todoitem{最后重写 abstract 和 conclusion 中关于数值结果的两三句话。}
\end{itemize}

\subsection{暂时不做}

\begin{itemize}
    \item 不证明完整 fully discrete AP theorem。
    \item 不扩展新的理论假设。
    \item 不新增过多复杂模型。
    \item 不把 density-compatible variant 写成第二个主算法。
    \item 不追求 direct density 与 WKB 在尖峰逐点形状上的完全一致。
\end{itemize}

\section{执行原则}

这几天的核心原则是控制范围。理论部分只收紧，不扩展。全离散部分只说明策略，不强行证明。数值部分先定实验逻辑，不急着写结果。新增实验只服务一个目标，即说明 WKB 粗网格重构能够达到 direct fine trait grid 的有效分辨效果。

因此，本阶段的最终目标可以概括为
\begin{center}
    \fbox{\large 非数值部分完全冻结，数值部分只等图。}
\end{center}

\end{document}
